\chapter{Comparison with web frameworks}
\label{ch:web}

\begin{aside}
\maya{}'s declarative model borrows from React, Svelte, Vue,
SolidJS. The borrowed ideas, the deliberate differences.
\end{aside}

\section{React}

\textbf{Borrowed:} declarative components, virtual DOM
analog, hooks model.

\textbf{Different:} no JSX (\cpp{} doesn't allow);
explicit lifetimes via owners; element tree rebuilt every
frame (no virtual DOM diff at element level).

\section{Svelte}

\textbf{Borrowed:} signals as the unit of reactivity.

\textbf{Different:} no compiler; everything at runtime.

\section{Vue}

\textbf{Borrowed:} reactivity API (refs vs reactive).

\textbf{Different:} no templates; only \cpp{}.

\section{SolidJS}

\textbf{Borrowed:} fine-grained reactivity. The closest
inspiration.

\textbf{Different:} \cpp{} type system; no JSX.

\section{Flutter}

\textbf{Borrowed:} build method; widget tree.

\textbf{Different:} text grid not pixel; less animation
emphasis.

\section{Why a new library}

The TUI niche is small, and existing TUI libs are largely
imperative. The web's declarative model fits this medium
beautifully. \maya{} is the bridge.

\section{Recap}

\begin{recap}
\begin{itemize}[leftmargin=*]
  \item Borrowed: signals, components, hooks-like
    lifecycle.
  \item Differences: \cpp{}, no JSX, terminal-native.
  \item SolidJS is the closest cousin.
\end{itemize}
\end{recap}

\section*{Workshop}

\begin{enumerate}[leftmargin=*]
  \item Compare a React component to its \maya{}
    equivalent line-for-line.
  \item Identify which features port and which don't.
\end{enumerate}
